We design four domain-specific heuristics to guide A* search toward reliable, low-risk, and realistically executable arbitrage routes. Each heuristic captures a different dimension of execution feasibility. Rather than detecting theoretical profit first and evaluating risk afterward, our approach embeds execution awareness directly into the search process.

\subsection{Design Principles for Domain-Specific Heuristics}

The heuristics are constructed according to three design principles:

\begin{enumerate}
    \item \textbf{Execution Awareness:} Heuristics must model real-world trading frictions such as liquidity limits, slippage, and transfer latency.
    \item \textbf{Early Pruning:} Infeasible or high-risk paths should be deprioritized before deep exploration.
    \item \textbf{Computational Efficiency:} Heuristic computation must remain lightweight relative to node expansion cost.
\end{enumerate}

Each heuristic modifies the evaluation function
\[
f(n) = g(n) + h(n),
\] % this can be inline
where $g(n)$ is accumulated cost and $h(n)$ encodes domain knowledge about future execution feasibility.

\subsection{$h_1$ -- Liquidity Heuristic}

The liquidity heuristic models the risk that low-volume markets cannot support the required trade size without excessive execution impact.

Let
\begin{align}
\textit{vol\_per\_sec}
&=
\frac{\textit{24h\_quote\_volume}}{24 \times 3600}
\\
\textit{liquidity\_ratio}
&=
\frac{\textit{vol\_per\_sec} \cdot T_{\textit{remain}}}
{\textit{order\_size\_usd}}.
\end{align}

This ratio estimates whether expected market flow during the remaining arbitrage window can support the intended order size.

We normalize this quantity into a bounded score:

\[
\textit{score}_{liq} =
\textit{clamp}\!\left(
\frac{\log_{10}(\textit{liquidity\_ratio}) + 1}{2},
0,
1
\right).
\]

The heuristic cost becomes:

\[
h_1(n) = \lambda_{liq} \big(1 - \textit{score}_{liq}(n)\big).
\]

\textbf{Characteristics:}
\begin{itemize}
    \item Time-aware liquidity modeling
    \item Accounts for order size relative to trading activity
    \item Favors deep, actively traded markets
\end{itemize}

This heuristic is well suited for conservative strategies prioritizing reliable execution over aggressive profit maximization.

\subsection{$h_2$ -- Slippage Heuristic}

The slippage heuristic directly models expected execution price impact using real-time order-book depth.

Let $P_{mid}$ denote the mid-price and $\textit{VWAP}(Q)$ the volume-weighted average price required to fill order size $Q$. Slippage in basis points is defined as:

\[
\textit{slippage}_{bps} =
\frac{\textit{VWAP}(Q) - P_{mid}}{P_{mid}} \times 10^{4}.
\]

The heuristic is defined as:

\[
h_2(n) =
\lambda_{slip} \cdot 
\max \left\{ 0, \textit{slippage}_{bps} - \theta \right\},
\]

where $\theta$ is a configurable slippage tolerance (default: 10 bps).

\textbf{Characteristics:}
\begin{itemize}
    \item Uses real-time order book data
    \item Side-aware (buy vs. sell impact)
    \item Directly penalizes large price impact
\end{itemize}

This heuristic performs particularly well when order-book depth is the dominant constraint.

\subsection{$h_3$ -- Parallel Execution Heuristic}

The parallel-execution heuristic is a meta-strategy that runs multiple A* searches concurrently from diverse starting nodes. Each parallel instance uses a base heuristic (typically $h_1$ or $h_2$), and the best discovered solution is selected.

Because $h_3$ operates at the search-control level rather than as a scalar node estimate, it does not admit a closed-form cost expression. Instead, it modifies exploration dynamics by expanding multiple regions of the state space simultaneously.

\textbf{Characteristics:}
\begin{itemize}
    \item Diversifies search across starting exchanges
    \item Mitigates local optima bias
    \item Higher computational cost
\end{itemize}

This approach is particularly useful when profitable starting points are not known in advance.

\subsection{$h_4$ -- Chain Congestion and Exchange Reliability Heuristic}

The $h_4$ heuristic incorporates transfer latency and operational reliability into the search.

Let $t_{min}$ denote the shortest outgoing transfer time from the current node, and $T_{remain}$ the remaining arbitrage window. The chain reliability score is:

\[
\textit{score}_{chain} =
1 - \frac{t_{min}}{T_{remain}}.
\]

The corresponding chain penalty is:

\[
h_{chain}(n) =
\lambda_{chain} \big(1 - \textit{score}_{chain}(n)\big).
\]

If each exchange is assigned a reliability score $s(n) \in [0,1]$, we define:

\[
h_{exchange}(n) =
\lambda_{exchange} \big(1 - s(n)\big).
\]

The combined heuristic is:

\[
h_4(n) = h_{chain}(n) + h_{exchange}(n).
\]

\textbf{Characteristics:}
\begin{itemize}
    \item Models blockchain congestion risk
    \item Incorporates exchange reliability
    \item Prioritizes fast and stable execution paths
\end{itemize}

In our experiments, this heuristic produced the fastest and most robust arbitrage routes, particularly in short-window environments.